% © 2025 Ing. David Jaroš — CC BY-NC-ND 4.0
%
% This work is licensed under a Creative Commons Attribution-NonCommercial-NoDerivatives
% 4.0 International License (CC BY-NC-ND 4.0).
%
% License History: Earlier drafts (up to v0.3) were released under CC BY 4.0.
% From v0.4 onward, all material is released under CC BY-NC-ND 4.0 to protect
% the integrity of the theoretical work during ongoing academic development.
%
% See LICENSE.md for full license text.

\documentclass[12pt]{article}
\usepackage[a4paper, margin=2.5cm]{geometry}
\usepackage{amsmath, amssymb, amsthm}
\usepackage{booktabs}
\usepackage{longtable}
\usepackage{hyperref}
\usepackage{tcolorbox}

\newtheorem{theorem}{Theorem}[section]
\newtheorem{lemma}[theorem]{Lemma}
\newtheorem{corollary}[theorem]{Corollary}
\newtheorem{proposition}[theorem]{Proposition}
\theoremstyle{definition}
\newtheorem{definition}[theorem]{Definition}
\theoremstyle{remark}
\newtheorem{remark}[theorem]{Remark}

\title{\textbf{The GR Limit of Unified Biquaternion Theory:\\
       Combined Variation, Admissible Sector, and\\
       the Role of the $\Theta$ Field}}
\author{Ing.\ David Jaroš}
\date{March 2026}

\begin{document}
\maketitle
\tableofcontents
\bigskip

\begin{abstract}
This document provides the single, canonical derivation of how Unified
Biquaternion Theory (UBT) reproduces Einstein's field equations
$G_{\mu\nu} = 8\pi G\,T_{\mu\nu}$ in the general-relativistic sector.
Three themes are unified here:

\begin{enumerate}
  \item \textbf{The $\Theta$ field as a generating field.}  The biquaternionic
    field $\Theta(q,\tau)$ is a potential-like object from which both spacetime
    geometry (the metric $g_{\mu\nu}$) and gauge fields emerge.

  \item \textbf{Combined variation vs.\ termwise separation.}  Varying the
    pure-$\Theta$ action yields a single coupled Euler--Lagrange condition
    $\mathcal{E}_\Theta + J^{*}\mathcal{E}_g = 0$, not two independently
    vanishing terms.  Einstein's equations are extracted from this combined
    condition under explicit regularity assumptions.

  \item \textbf{The admissible GR sector.}  GR is reproduced \emph{exactly on
    the admissible class} $\mathcal{A}_{\mathrm{UBT}}$ defined by conditions
    C1--C6.  Outside that sector, UBT may admit additional non-GR solutions,
    representing an extension of Einstein gravity rather than a contradiction.
\end{enumerate}

\medskip\noindent
\textbf{Correct language for GR recovery:}
\begin{itemize}
  \item[\checkmark] ``GR limit reproduced on the admissible sector'' — accurate.
  \item[\checkmark] ``UBT generalises GR'' — accurate.
  \item[$\times$] ``GR proved (universally)'' — overstated; recovery is conditional.
  \item[$\times$] ``UBT contradicts GR'' — incorrect.
\end{itemize}
\end{abstract}

% =========================================================================
\section{The $\Theta$ Field as a Generating Field}
\label{sec:theta_role}

\subsection{Definition and structure}

The fundamental object of UBT is the biquaternionic field
\begin{equation}
  \Theta: M^4 \times \mathcal{C} \;\longrightarrow\; \mathcal{B},
\end{equation}
where $M^4$ is four-dimensional spacetime, $\mathcal{C}$ is the complex-time
fibre (parameterised by $\tau = t + i\psi$), and $\mathcal{B} = \mathbb{C}\otimes\mathbb{H}$
is the biquaternion algebra.  The field $\Theta$ is \emph{not} itself the
metric or a gauge connection; it is a more fundamental, potential-like object
from which these geometric structures emerge.

\subsection{Emergent metric}

The spacetime metric is defined as the real part of the normalised outer
product of $\Theta$ with its quaternionic conjugate:
\begin{equation}
  g_{\mu\nu}[\Theta](x)
  = \operatorname{Re}\!\left[
      \frac{\partial_\mu\Theta\cdot\partial_\nu\Theta^\dagger}
           {\mathcal{N}[\Theta]}
    \right],
  \label{eq:emergent_metric}
\end{equation}
where $\mathcal{N}[\Theta] = \operatorname{Re}[\partial^\mu\Theta\cdot\partial_\mu\Theta^\dagger]$
is the normalisation scalar.  This formula is the \emph{derivative-based
metric}, proved equivalent to the tetrad-based metric under AXIOM~B and the
metric bridge theorem (see \texttt{GR\_closure/step1\_metric\_bridge.tex}).

\subsection{Emergent gauge fields}

Beyond the metric, the full biquaternion structure of $\Theta$ encodes gauge
connections.  The imaginary-part components of $\Theta$ generate the $\mathrm{U}(1)$,
$\mathrm{SU}(2)_L$, and $\mathrm{SU}(3)_c$ gauge fields of the Standard Model
through the automorphism structure of $\mathbb{C}\otimes\mathbb{H}$ (see
\texttt{consolidation\_project/appendix\_E\_SM\_QCD\_embedding.tex} and
\texttt{appendix\_E2\_SM\_geometry.tex}).

\subsection{Layer architecture}

The generating role of $\Theta$ gives rise to the following conceptual
hierarchy:

\begin{center}
\begin{tabular}{lll}
\toprule
\textbf{Layer} & \textbf{Content} & \textbf{Status} \\
\midrule
Layer 0 & Axioms, $\mathcal{B}$-algebra, $\Theta$ definition,
          complex time $\tau = t+i\psi$ & Locked \\
Layer 1 & Physical sector: GR, Standard Model gauge fields,
          emergent metric & Active derivation \\
Layer 2 & Informational/coding structures: winding numbers,
          channel selection, discrete structures & Hypothesis \\
\bottomrule
\end{tabular}
\end{center}

Layer~0 provides the axioms and mathematical structure from which all other
layers are derived.  Layer~1 contains the physical predictions (GR, gauge
fields) obtained by projecting the Layer~0 structure onto observable degrees
of freedom.  Layer~2 contains informational and coding structures that select
specific physical channels within the Layer~1 solution space.

% =========================================================================
\section{The Total Action and Its Variation}
\label{sec:action_variation}

\subsection{Setup: total action}

Let the total action be
\begin{equation}
  S_{\text{total}}[g, \Theta]
  = S_{\text{EH}}[g] + S_\Theta[g, \Theta],
  \label{eq:total_action}
\end{equation}
where $S_{\text{EH}}[g] = \frac{1}{16\pi G}\int\!\sqrt{-g}\,R\,d^4x$ is the
Einstein--Hilbert action and $S_\Theta[g,\Theta]$ contains all $\Theta$-dependent
terms (kinetic term, coupling to matter, etc.).

\subsection{Euler--Lagrange operators}
\label{subsec:el_ops}

\begin{definition}[Euler--Lagrange operators]
\label{def:el_operators}
\begin{align}
  \mathcal{E}_\Theta[\Theta, g]
  &:= \frac{\delta S_\Theta}{\delta\Theta}\bigg|_{g \text{ fixed}},
  \label{eq:E_Theta}\\[4pt]
  \mathcal{E}_g[g, \Theta]
  &:= \frac{\delta S_{\text{total}}}{\delta g^{\mu\nu}}\bigg|_{\Theta \text{ fixed}}
   = \frac{1}{16\pi G}\!\left(G_{\mu\nu} - 8\pi G\, T_{\mu\nu}[\Theta,g]\right),
  \label{eq:E_g}
\end{align}
and the \emph{induced variation map}
\begin{equation}
  J\,\delta\Theta := \frac{\delta g_{\mu\nu}[\Theta]}{\delta\Theta}\,\delta\Theta.
  \label{eq:J_map}
\end{equation}
The formal $L^2$-adjoint of $J$ is denoted $J^*$.
\end{definition}

\subsection{Pure-$\Theta$ functional}

Substituting the emergent metric $g_{\mu\nu}[\Theta]$ from
Eq.~\eqref{eq:emergent_metric} into $S_{\text{total}}$, one obtains the
pure-$\Theta$ functional
\begin{equation}
  \hat{S}[\Theta] := S_{\text{total}}\bigl[g[\Theta],\,\Theta\bigr].
  \label{eq:S_hat}
\end{equation}

Applying the chain rule to $\delta\hat{S}/\delta\Theta$:
\begin{equation}
  \frac{\delta\hat{S}}{\delta\Theta}\,\delta\Theta
  = \mathcal{E}_\Theta\,\delta\Theta
  + \mathcal{E}_g^{\mu\nu}\,J\,\delta\Theta,
  \label{eq:variation_chain}
\end{equation}
so stationarity $\delta\hat{S} = 0$ for all admissible $\delta\Theta$ reads:

% =========================================================================
\section{Combined Variation and the Admissible GR Sector}
\label{sec:combined_variation}

\begin{tcolorbox}[colback=blue!5!white,colframe=blue!75!black,
                  title=Fundamental Variational Equation of UBT]
The Euler--Lagrange equation obtained from the pure-$\Theta$ action is the
\textbf{combined condition}
\begin{equation}
  \mathcal{E}_\Theta + J^{*}\mathcal{E}_g = 0.
  \tag{$\star$}
  \label{eq:combined_EL}
\end{equation}
This is a \textbf{single} coupled equation, \emph{not} two equations
$\mathcal{E}_\Theta = 0$ and $\mathcal{E}_g = 0$ imposed independently.
\end{tcolorbox}

\subsection{Why termwise separation does not follow automatically}

\begin{remark}[No automatic separation]
\label{rem:no_separation}
Even though Eq.~\eqref{eq:combined_EL} holds for all $\delta\Theta$, one
\emph{cannot} conclude that $\mathcal{E}_\Theta = 0$ and $\mathcal{E}_g = 0$
independently.  The two terms both depend on the \emph{same} variation
$\delta\Theta$; they are not independently varied.  The correct statement is
only that the linear combination $\mathcal{E}_\Theta + J^{*}\mathcal{E}_g$
vanishes.
\end{remark}

\begin{tcolorbox}[colback=yellow!10!white,colframe=orange!80!black,
                  title=Warning: Do Not Assume Independent Termwise Vanishing]
It is mathematically \emph{incorrect} to argue: ``since the variation holds for
arbitrary $\delta\Theta$, each term must vanish separately.''  Separate vanishing
requires $\delta\Theta$ to vary each term independently, which is generally
impossible when both terms depend on the same $\delta\Theta$.
\end{tcolorbox}

\subsection{When does $\mathcal{E}_g = 0$ (Einstein's equations) follow?}
\label{subsec:when_sep}

Separate vanishing $\mathcal{E}_g = 0$ — i.e.\ Einstein's equations
$G_{\mu\nu} = 8\pi G\,T_{\mu\nu}$ — follows from Eq.~\eqref{eq:combined_EL}
under any one of the following conditions:

\begin{enumerate}
  \item \textbf{Surjectivity / rank condition on $J$:} If $J$ is locally
    surjective (full column rank), any metric variation $\delta g^{\mu\nu}$
    can be written as $J\,\delta\Theta$, so
    $\langle\mathcal{E}_g, J\,\delta\Theta\rangle = 0$ for all $\delta\Theta$
    implies $\mathcal{E}_g = 0$.

  \item \textbf{Gauge-reduced local injectivity:} After quotienting by metric-invisible
    gauge directions in $\Theta$-space, the map $\Theta\mapsto g[\Theta]$ is
    locally invertible, so $\Theta$-variations span the full metric variation
    space.

  \item \textbf{Projection to metric observables:} If all non-metric $\Theta$
    modes are either gauge redundancies or frozen (suppressed at classical
    scales), the contracted form of Eq.~\eqref{eq:combined_EL} reduces to
    $\mathcal{E}_g = 0$.

  \item \textbf{On-shell $\mathcal{E}_\Theta = 0$:} If $\Theta$ satisfies its
    own field equation separately, Eq.~\eqref{eq:combined_EL} reduces to
    $J^{*}\mathcal{E}_g = 0$, and condition~(1) then implies $\mathcal{E}_g = 0$.
\end{enumerate}

\begin{theorem}[GR recovery on the admissible sector]
\label{thm:gr_recovery}
Suppose $\Theta \in \mathcal{A}_{\mathrm{UBT}}$ (Definition~\ref{def:admissible}).
Then the combined condition~\eqref{eq:combined_EL} implies
\begin{equation}
  G_{\mu\nu} = 8\pi G\, T_{\mu\nu}[\Theta, g[\Theta]].
  \label{eq:einstein}
\end{equation}
\end{theorem}

\begin{proof}[Proof sketch]
Under the rank condition (conditions~1 or~2), any metric variation
$\delta g^{\mu\nu}$ is of the form $J\,\delta\Theta$ for some $\delta\Theta$.
Therefore $\langle\mathcal{E}_g,\,J\,\delta\Theta\rangle = 0$ for all
$\delta\Theta$ implies $\langle J^*\mathcal{E}_g,\,\delta\Theta\rangle = 0$
for all $\delta\Theta$.  Surjectivity of $J$ then gives $\mathcal{E}_g = 0$.
\hfill$\square$
\end{proof}

\subsection{Outside the admissible sector}

In the general case, outside $\mathcal{A}_{\mathrm{UBT}}$:
\begin{itemize}
  \item The kernel of $J$ may be non-trivial: some $\delta\Theta$ directions
    do not change the metric.  These ``metric-invisible'' directions are
    genuine UBT degrees of freedom absent from GR.
  \item The combined condition may be satisfied by configurations where
    $\mathcal{E}_g \neq 0$ (non-GR solutions of UBT).
  \item This is \emph{not a defect}: UBT contains GR as a proper subsector
    while admitting additional dynamics.
\end{itemize}

% =========================================================================
\section{The Admissible Configuration Space}
\label{sec:admissible}

\begin{definition}[Admissible class $\mathcal{A}_{\mathrm{UBT}}$]
\label{def:admissible}
A $\Theta$ configuration belongs to $\mathcal{A}_{\mathrm{UBT}}$ if and only
if conditions C1--C6 (Table~\ref{tab:conditions}) hold at every point of $M^4$.
\end{definition}

\begin{remark}
The admissible class is non-empty: flat Minkowski space ($\Theta = \Theta_0$
constant) belongs to $\mathcal{A}_{\mathrm{UBT}}$, as do all linearised GR
perturbations studied in the GR recovery proof.  Whether every classical GR
solution lifts to $\mathcal{A}_{\mathrm{UBT}}$ at the non-perturbative level
is an open problem.
\end{remark}

% -------------------------------------------------------------------------
\subsection{Gauge redundancies}

The map $\Theta \mapsto g[\Theta]$ is not injective: multiple $\Theta$
configurations may induce the same metric.  The kernel of $J$ consists of
``metric-invisible'' $\Theta$ directions.  Before asking whether the map is
invertible, one must quotient by this gauge redundancy.

\begin{definition}[Gauge-reduced local injectivity]
\label{def:gr_injectivity}
The map $\Theta \mapsto g[\Theta]$ is \emph{gauge-reduced locally injective}
at $\Theta_0$ if the induced map on the quotient $T_{\Theta_0}\mathcal{F}/\ker J$
is injective: for all $\delta\Theta \notin \ker J$, $J\,\delta\Theta \neq 0$.
\end{definition}

% -------------------------------------------------------------------------
\subsection{Classical limit and decoupling assumptions}

The GR recovery result is a \emph{classical} statement.  It assumes:
\begin{enumerate}
  \item $\hbar \to 0$: the path integral is dominated by the saddle point.
  \item The Lorentzian sector dominates: imaginary-time corrections are
    suppressed by $\psi/t \ll 1$.
  \item The biquaternionic field equation reduces to its real projection
    when quantum and imaginary-time effects are suppressed.
\end{enumerate}

% =========================================================================
\section{Condition Table and Proof Status}
\label{sec:conditions}

\begin{tcolorbox}[colback=gray!5!white,colframe=gray!60!black,
                  title=Mandatory Caveat]
\textbf{GR reproduction is exact only on the sector satisfying conditions
C1--C6 below.  Outside that sector, UBT may admit additional non-GR
solutions.  No claim is made that GR is universally derived from UBT without
these conditions.}
\end{tcolorbox}

\bigskip

\begin{longtable}{@{}p{3.2cm}p{3.4cm}p{3.5cm}p{1.8cm}p{2.0cm}@{}}
\caption{Conditions for exact GR recovery in UBT.\label{tab:conditions}}\\
\toprule
\textbf{Condition} &
\textbf{Physical meaning} &
\textbf{Required for} &
\textbf{Proved} &
\textbf{Status}
\\
\midrule
\endfirsthead
\multicolumn{5}{c}{\textit{(continued)}}\\
\toprule
\textbf{Condition} &
\textbf{Physical meaning} &
\textbf{Required for} &
\textbf{Proved} &
\textbf{Status}
\\
\midrule
\endhead
%
C1.\ Non-degeneracy of $\Theta \to g$ map
& Metric non-degenerate; no metric-invisible $\Theta$ directions except gauge.
& Gauge-reduced injectivity; allows recovering $\mathcal{E}_g = 0$.
& Perturbative.
& Open non-perturbatively.
\\[4pt]
%
C2.\ Gauge-reduced uniqueness
& Gauge-equivalent $\Theta$ configurations yield the same metric.
& Uniqueness of emergent metric on physical states.
& Plausible.
& Gauge structure not fully classified.
\\[4pt]
%
C3.\ Regularity of induced metric
& $g_{\mu\nu}[\Theta]$ is smooth, Lorentzian $(-,+,+,+)$, satisfies metric postulates.
& Well-defined curvature and Bianchi identity.
& \textbf{Proved}: signature theorem (step3), non-degeneracy for $\mathcal{A}_{\mathrm{UBT}}$.
& —
\\[4pt]
%
C4.\ Existence of classical metric sector
& Every GR solution $g_{\mu\nu}$ can be lifted to some $\Theta \in \mathcal{A}_{\mathrm{UBT}}$.
& GR is embeddable in UBT.
& Proved for flat + linearised; non-linear: open.
& Assumed at non-linear level.
\\[4pt]
%
C5.\ Decoupling of extra $\Theta$ modes
& Non-metric $\Theta$ degrees of freedom (imaginary part, $\psi$-modes) are dynamically suppressed.
& Reduces combined EL equation to Einstein equation alone.
& Not proved in general.
& Assumed for classical gravity.
\\[4pt]
%
C6.\ Bianchi compatibility
& $\nabla^\mu G_{\mu\nu} = 0$ and $\nabla^\mu T_{\mu\nu} = 0$ after projection.
& Consistency of GR.
& \textbf{Proved}: contracted Bianchi identity under A1--A2.
& —
\\
\bottomrule
\end{longtable}

\begin{remark}[Summary of proof status]
Conditions C3 and C6 are \textbf{fully proved} within UBT.
Conditions C1 and C4 are \textbf{proved in the linearised (perturbative) regime}
and remain open non-perturbatively.
Conditions C2 and C5 are \textbf{physically motivated assumptions} consistent
with UBT structure; formal proofs are listed as open problems in
\texttt{AUDITS/followup\_gap\_backlog\_2026\_03.md}.
\end{remark}

% =========================================================================
\section{Three Levels of GR Recovery}
\label{sec:three_levels}

Depending on the strength of conditions imposed, three levels can be
distinguished (see also \texttt{reports/gr\_recovery\_levels.md}):

\begin{enumerate}
  \item \textbf{Exact projected recovery (Level 1):} A clean projection theorem
    maps the $\Theta$ equation directly to $G_{\mu\nu} = 8\pi G\,T_{\mu\nu}$
    without additional assumptions.
    \emph{Status: open problem; requires full rank condition (C1) non-perturbatively.}

  \item \textbf{Constrained sector recovery (Level 2):} On $\mathcal{A}_{\mathrm{UBT}}$,
    the induced dynamics reproduces Einstein gravity.
    \emph{Status: current best-established result; conditions stated explicitly
    in this document.}

  \item \textbf{Low-energy effective limit (Level 3):} After freezing
    additional $\Theta$ degrees of freedom that decouple at macroscopic scales,
    the theory reduces to GR in the infrared.
    \emph{Status: plausible; requires analysis of the $\Theta$-mode spectrum.}
\end{enumerate}

\textbf{The current standing of UBT is Level~2.}

% =========================================================================
\section{Why UBT $\supseteq$ GR Is Not a Contradiction}
\label{sec:no_contradiction}

\begin{tcolorbox}[colback=green!8!white,colframe=green!60!black,
                  title=Consistency Statement]
There is \textbf{no contradiction} between UBT and GR.  UBT reproduces GR
on the admissible sector $\mathcal{A}_{\mathrm{UBT}}$ while containing
additional non-GR degrees of freedom beyond that sector.  The absence of
automatic termwise separation is a \emph{feature}: it indicates UBT is
\emph{richer} than GR, not inconsistent with it.
\end{tcolorbox}

Standard analogies from physics illustrate this structure:

\begin{itemize}
  \item \textbf{Tetrad/Palatini formulation of GR:} Varying the tetrad and
    connection yields Einstein's equations plus a torsion-free condition;
    neither condition follows from the other without additional structure.
    UBT is analogous: $\Theta$ is fundamental and the metric is derived.

  \item \textbf{EFT / integrating out heavy fields:} In effective field theory,
    the low-energy action does not satisfy the UV field equations termwise;
    agreement holds only on-shell or in a specific limit.  GR recovery from
    UBT is structurally of this type.

  \item \textbf{Precedent from other fundamental theories:} String theory
    recovers GR as a low-energy limit; loop quantum gravity recovers GR in
    the classical limit; both reproduce GR on a sector, not universally.
\end{itemize}

\medskip\noindent
\textbf{Correct and incorrect formulations:}
\begin{itemize}
  \item[\checkmark] ``GR limit reproduced on admissible sector $\mathcal{A}_{\mathrm{UBT}}$''
  \item[\checkmark] ``UBT generalises GR''
  \item[\checkmark] ``UBT extends GR with additional degrees of freedom''
  \item[\checkmark] ``GR is a sector of UBT: $\mathrm{GR}\subset\mathcal{A}_{\mathrm{UBT}}\subset\mathrm{UBT}$''
  \item[$\times$] ``GR proved (universally)''
  \item[$\times$] ``UBT contradicts GR''
  \item[$\times$] ``Absence of termwise separation weakens UBT fatally''
\end{itemize}

% =========================================================================
\section{Conclusion}
\label{sec:conclusion}

The fundamental equation of UBT in the $\Theta$-only formulation is the
combined variational condition
\begin{equation}
  \mathcal{E}_\Theta + J^{*}\mathcal{E}_g = 0.
  \label{eq:combined_final}
\end{equation}
Einstein's field equations $G_{\mu\nu} = 8\pi G\,T_{\mu\nu}$ are recovered
from this condition on the admissible class $\mathcal{A}_{\mathrm{UBT}}$
(Level~2 recovery).  The mathematical structure is summarised by:
\begin{equation}
  \mathrm{GR} \subset \mathcal{A}_{\mathrm{UBT}} \subset \mathrm{UBT}.
\end{equation}

The $\Theta$ field plays the role of a generating field: spacetime geometry
and gauge fields emerge from it through the projection maps described in
Section~\ref{sec:theta_role}.  The three-layer architecture (Layer~0:
axioms and algebra; Layer~1: physical sector; Layer~2: informational/coding
structures) organises this derivation hierarchy.

\medskip\noindent
\textbf{Source documents merged into this file:}
\begin{itemize}
  \item \texttt{consolidation\_project/GR\_closure/theta\_vs\_metric\_variation\_note.tex}
    — combined variation and termwise separation analysis.
  \item \texttt{consolidation\_project/GR\_closure/gr\_sector\_conditions.tex}
    — explicit conditions C1--C6 and condition table.
\end{itemize}

\medskip\noindent
\textbf{Related documents:}
\begin{itemize}
  \item \texttt{consolidation\_project/GR\_closure/GR\_chain\_summary.tex}
    — step-by-step GR derivation chain.
  \item \texttt{reports/gr\_recovery\_levels.md} — summary of proof levels.
  \item \texttt{docs/ubt\_gr\_relationship.md} — non-technical overview.
  \item \texttt{core/AXIOMS.md} — Layer~0 axioms.
\end{itemize}

\end{document}
